\documentclass[11pt,a4paper]{article}
\usepackage[T1]{fontenc}
\usepackage[utf8]{inputenc}
\usepackage{lmodern,amsmath,amssymb}
\usepackage[margin=25mm]{geometry}
\usepackage[hidelinks]{hyperref}
\hypersetup{pdftitle={The probabilistic Planck gap: $F\tau(L+\tau P/2m)\ge\hbar\arcsin(1-2\epsilon)$, and why it},pdfauthor={navstokgap research working notes}}
\newcommand{\tightlist}{\setlength{\itemsep}{0pt}\setlength{\parskip}{0pt}}
\setlength{\emergencystretch}{3em}
\setcounter{secnumdepth}{0}
\begin{document}
\hypertarget{the-probabilistic-planck-gap-ftaultau-p2mgehbararcsin1-2epsilon-and-why-it-is-resource-relative}{%
\section{\texorpdfstring{The probabilistic Planck gap:
\(F\tau(L+\tau P/2m)\ge\hbar\arcsin(1-2\epsilon)\), and why it is
resource-relative}{The probabilistic Planck gap: F\textbackslash tau(L+\textbackslash tau P/2m)\textbackslash ge\textbackslash hbar\textbackslash arcsin(1-2\textbackslash epsilon), and why it is resource-relative}}\label{the-probabilistic-planck-gap-ftaultau-p2mgehbararcsin1-2epsilon-and-why-it-is-resource-relative}}

\textbf{Sharpening, 2026-09-23.} Every \(1-2\epsilon\) in the protocol
statements below improves to
\(\arcsin(1-2\epsilon)=\arccos(2\sqrt{\epsilon(1-\epsilon)})\), the
single-shot constant of Theorem 1, by running Theorem 2's hybrid chain
with the Bures angle in place of total variation (Theorem P of the
\href{record-distance-path-length.md}{path-length note}); the proof of
Theorem 2 as written gives the weaker constant. That note also gives the
kick form \(\int_0^\tau|f(t)|L(t)\,dt\ge\hbar\arcsin(1-2\epsilon)\),
with \(L(t)\) the position spread at time \(t\), from which Theorem 2
follows.

\textbf{Correction, 2026-09-22.} The first version stated Corollary 3 at
fixed area \(A=LP\) with the aperture's shape free. Its proof minimized
\(2mL+\tau A/L\) over \(L\), which yields the largest lower bound,
attained only at the balanced aperture \(L=\tau P/2m\). For an aperture
of aspect \(r=L/\sqrt{\tau A/2m}\) the floor is
\((1-2\epsilon)^2\hbar^2/(A(r+r^{-1})^2)\), and a squeezed Gaussian of
area \(\hbar/2\) and large \(r\) drives it to zero. Theorems 1 and 2,
stated in \(L\) and \(P\) separately, are unaffected; Corollary 3, the
lead and Sections 5--6 now carry the balanced form.

The worst-case theorem of the \href{planck-gap-derivation.md}{derivation
note} has no quantum instance. Its probabilistic replacement does, and
it decides the universality question that the worst-case version left
open. Let the apparatus confine the particle's position spread to \(L\)
and its momentum spread to \(P\), and set \(A=LP\), the phase-space area
available to a preparation; the uncertainty relation forces
\(A\ge\hbar/2\). Then for \textbf{every} protocol of quantum
instruments, of any number, at any times, with any intermediate outcomes
and any final measurement, distinguishing the falling from the inertial
hypothesis with error probability at most \(\epsilon\) requires

\[\boxed{\frac{F\tau L}{\hbar}+\frac{F\tau^2P}{2m\hbar}\ge\arcsin(1-2\epsilon),
\qquad\text{hence}\qquad
\tau\Delta E=\frac{F^2\tau^3}{2m}\ge\frac{\arcsin^2(1-2\epsilon)\,\hbar^2}{A\,(r+r^{-1})^2},}\]

with \(r=L/\sqrt{\tau A/2m}\) the aperture's aspect measured against the
duration (Theorem 2, Corollary 3). In Newton's quantities the first form
reads \(J L+sP\ge\hbar\arcsin(1-2\epsilon)\): the impulse \(J=F\tau\) of
Proposition I weighed against the position aperture, plus the sagitta
\(s=F\tau^2/2m\) of Lemma X weighed against the momentum aperture. At
the balanced aperture \(r=1\), where the two terms are equal, and at the
minimal area \(A=\hbar/2\), the floor is
\(\tau\Delta E\ge\frac12\arcsin^2(1-2\epsilon)\,\hbar\), the same
combination \(F^2\tau^3/m\) that Newton's geometry sends to zero and
that the classical theorem bounds by the marks' phase-space cost. The
single-shot bound (Theorem 1) and the protocol bound share the constant
\(\arcsin(1-2\epsilon)=\arccos(2\sqrt{\epsilon(1-\epsilon)})\), which is
\(\pi/2\) at \(\epsilon=0\).

\textbf{The universality question is decided, and the answer is negative
without the resource bound.} For every \(\tau\Delta E>0\) there is a
preparation and a measurement distinguishing the hypotheses with error
probability \(O(\tau\Delta E/\hbar)\), built from two narrow packets
separated by \(d=\pi\hbar/(F\tau)\) (Proposition 4). The separation
grows without bound as the force shrinks, saturating Theorem 2, so the
theorem is tight in \(L\) and the failure of universality is exactly the
failure of confinement. Protocol-universality over instruments survives
in full; universality over preparations fails. Ozawa's measurements
breaking the standard quantum limit for free-mass position belong to the
same phenomenon and are covered by the resource-relative statement.

Methods are standard: Mandelstam--Tamm for the single-shot bound, a
telescoping hybrid argument for the many-instrument case. The
contribution is the protocol-universal statement in the Galileo geometry
with explicit constants, and a counterexample that matches it.
Exploratory; no ledger promotion.

\hypertarget{the-two-hypotheses-differ-by-a-phase-space-displacement}{%
\subsection{1. The two hypotheses differ by a phase-space
displacement}\label{the-two-hypotheses-differ-by-a-phase-space-displacement}}

Transverse coordinate \(\hat y\), momentum \(\hat p\), mass \(m>0\), on
\([0,\tau]\). Hypothesis \(\mathrm I\) evolves with
\(H_{\mathrm I}=\hat p^2/2m\); hypothesis \(\mathrm F\) with
\(H_{\mathrm F}=\hat p^2/2m-F\hat y\). Write

\[D(a,b)=\exp\!\left[\frac{i}{\hbar}\left(b\hat y-a\hat p\right)\right],
\qquad D^\dagger\hat yD=\hat y+a,\qquad D^\dagger\hat pD=\hat p+b .\]

\textbf{Lemma 1.}
\(W(t)=U_{\mathrm I}(t)^\dagger U_{\mathrm F}(t) =e^{i\varphi(t)}D\!\left(\alpha(t),\beta(t)\right)\)
with

\[\alpha(t)=-\frac{Ft^2}{2m},\qquad \beta(t)=Ft,\]

and a real phase \(\varphi\).

\emph{Proof.} Under \(H_{\mathrm F}\) the Heisenberg operators are
\(\hat y(t)=\hat y+\hat pt/m+Ft^2/2m\) and \(\hat p(t)=\hat p+Ft\), so
\(U_{\mathrm F}(t)=e^{i\varphi}D(Ft^2/2m,Ft)\,U_{\mathrm I}(t)\) for
some phase. Free evolution conjugates a displacement to a displacement,
\(U_{\mathrm I}(t)^\dagger D(a,b)U_{\mathrm I}(t)=D(a-bt/m,b)\), and
\(Ft^2/2m-Ft\cdot t/m=-Ft^2/2m\). \(\square\)

The product of the two components is the programme's combination:

\[|\alpha(\tau)\beta(\tau)|=\frac{F\tau^2}{2m}\cdot F\tau
=\frac{F^2\tau^3}{2m}=\tau\Delta E=\frac{3F}{v}A_{\rm inertial,fall},\]

with \(\Delta E=F^2\tau^2/2m\) and the area of
\href{newton-insertion-action.md}{N01}. The displacement's position
component is the sagitta of Lemma X and its momentum component is the
accumulated impulse of Proposition I.

\textbf{Resource bound.} The apparatus is said to have \textbf{aperture}
\((L,P)\) when every state occurring in the experiment satisfies
\(\Delta_\psi\hat y\le L\) and \(\Delta_\psi\hat p\le P\), where
\(\Delta_\psi X=\|(X-\langle X\rangle_\psi)\psi\|\). Confinement to a
region of diameter \(2L\) and to momenta of modulus at most \(P\) gives
this for every state at once. Set \(A=LP\); since
\(\Delta\hat y\,\Delta\hat p\ge\hbar/2\) for any state that occurs,
\(A\ge\hbar/2\).

\hypertarget{single-shot}{%
\subsection{2. Single shot}\label{single-shot}}

\textbf{Theorem 1.} Let the preparation be a pure state \(\psi\) with
\(\Delta_\psi\hat y\le L\) and \(\Delta_\psi\hat p\le P\), let the two
hypotheses run undisturbed on \([0,\tau]\), and let any measurement on
the final state decide between them with equal priors and error
probability at most \(\epsilon<1/2\). Then

\[\frac{F\tau L}{\hbar}+\frac{F\tau^2P}{2m\hbar}\ \ge\
\theta(\epsilon):=\arccos\!\left(2\sqrt{\epsilon(1-\epsilon)}\right),
\qquad\theta(0)=\frac\pi2 .\]

\emph{Proof.} The two final states are \(U_{\mathrm I}\psi\) and
\(U_{\mathrm F}\psi=e^{i\varphi}U_{\mathrm I}D(\alpha,\beta)\psi\) by
Lemma 1 with \(\alpha=\alpha(\tau)\), \(\beta=\beta(\tau)\); a common
unitary and a global phase change no error probability, so the task is
to distinguish \(\psi\) from \(D(\alpha,\beta)\psi\). Helstrom's bound
for two pure states with equal priors gives
\(\epsilon\ge\frac12\bigl(1-\sqrt{1-|\langle\psi|D|\psi\rangle|^2}\bigr)\),
so \(|\langle\psi|D|\psi\rangle|\le2\sqrt{\epsilon(1-\epsilon)}\). Write
\(D=e^{-iG}\) with \(G=-(\beta\hat y-\alpha\hat p)/\hbar\). The
Mandelstam--Tamm bound in its geometric form states that the
Fubini--Study angle between \(\psi\) and \(e^{-isG}\psi\) grows at rate
at most \(\Delta_\psi G\), so
\(|\langle\psi|e^{-iG}|\psi\rangle|\ge\cos(\Delta_\psi G)\) whenever
\(\Delta_\psi G\le\pi/2\). Hence either
\(\Delta_\psi G>\pi/2\ge\theta(\epsilon)\), or
\(\cos(\Delta_\psi G)\le2\sqrt{\epsilon(1-\epsilon)}\) and again
\(\Delta_\psi G\ge\theta(\epsilon)\). Finally \(\Delta_\psi\) is a
seminorm on observables, being the norm of the centred operator applied
to \(\psi\), so

\[\Delta_\psi G\le\frac{|\beta|\Delta_\psi\hat y+|\alpha|\Delta_\psi\hat p}{\hbar}
\le\frac{F\tau L}{\hbar}+\frac{F\tau^2P}{2m\hbar}. \qquad\square\]

The bound uses no measurement model: it holds for every final
measurement, optimal or not, and it is a statement about the two states
alone.

\hypertarget{every-protocol-of-instruments}{%
\subsection{3. Every protocol of
instruments}\label{every-protocol-of-instruments}}

\textbf{Theorem 2.} Let instruments \(M_1,\dots,M_k\) act at times
\(0<t_1<\dots<t_k\le\tau\), identical under the two hypotheses, with any
Kraus operators, any quantum memory in the apparatus and arbitrary
dependence of later instruments on earlier outcomes, followed by any
final measurement. Let the apparatus have aperture \((L,P)\), in the
sense that the body's marginal of every conditional state arising in the
protocol satisfies the two spread bounds. If the full record decides
between the hypotheses with equal priors and error probability at most
\(\epsilon\), then

\[\frac{F\tau L}{\hbar}+\frac{F\tau^2P}{2m\hbar}\ \ge\ 1-2\epsilon .\]

\emph{Proof.} Pass to the interaction picture of \(U_{\mathrm I}\),
absorbing the free evolution into the instruments, so that under
\(\mathrm I\) the state is unchanged between instruments and under
\(\mathrm F\) it is acted on, between \(t_{r-1}\) and \(t_r\), by
\(V_r=W(t_r)W(t_{r-1})^\dagger=e^{i\varphi_r}D(\Delta\alpha_r,\Delta\beta_r)\)
with \(\Delta\alpha_r=\alpha(t_r)-\alpha(t_{r-1})\) and
\(\Delta\beta_r=\beta(t_r)-\beta(t_{r-1})\), using Lemma 1 and the
composition of displacements up to phase; set \(t_0=0\) and let
\(r=1,\dots,k+1\) with \(t_{k+1}=\tau\).

Let \(P_{\mathrm I}\) and \(P_{\mathrm F}\) be the distributions of the
full record. Define hybrid processes \(H_0,\dots,H_{k+1}\), where
\(H_s\) applies \(V_r\) for \(r\le s\) and the identity for \(r>s\);
then \(H_0=\mathrm I\) and \(H_{k+1}=\mathrm F\). Consecutive hybrids
differ by one insertion of \(V_s\) into an otherwise identical sequence,
so by the triangle inequality for total variation and the
data-processing inequality for the common remainder of the protocol,

\[\mathrm{TV}(P_{\mathrm I},P_{\mathrm F})\le\sum_{s}
\mathrm{TV}(P_{H_{s-1}},P_{H_s})\le\sum_s\ \sup_{\sigma}\ \tfrac12\bigl\|\sigma-V_s\sigma V_s^\dagger\bigr\|_1,\]

the supremum over the conditional states \(\sigma\) of body and
apparatus memory reachable at step \(s\), with \(V_s\) acting as
\(V_s\otimes\mathbf 1\). Purify \(\sigma\) to a vector \(\psi\). The
partial trace does not increase the trace norm, so
\(\tfrac12\|\sigma-V_s\sigma V_s^\dagger\|_1\le\sqrt{1-|\langle\psi|V_s\otimes\mathbf 1|\psi\rangle|^2}\),
and \(\Delta_\psi(G_s\otimes\mathbf 1)=\Delta_{\sigma_{\rm body}}G_s\)
because the first two moments of \(G_s\otimes\mathbf 1\) are those of
\(G_s\) in the body's marginal. The right side is at most
\(\sin(\Delta_\psi G_s)\le\Delta_\psi G_s\) when
\(\Delta_\psi G_s\le\pi/2\), by Mandelstam--Tamm as in Theorem 1, and at
most \(1\le\Delta_\psi G_s\) otherwise; in both cases the left side is
at most \(\Delta_{\sigma_{\rm body}}G_s\). Therefore

\[\mathrm{TV}(P_{\mathrm I},P_{\mathrm F})\le\sum_s
\frac{|\Delta\beta_s|L+|\Delta\alpha_s|P}{\hbar}
=\frac{L\sum_s|\Delta\beta_s|+P\sum_s|\Delta\alpha_s|}{\hbar}.\]

Both \(\alpha(t)=-Ft^2/2m\) and \(\beta(t)=Ft\) are monotone on
\([0,\tau]\), so the increments telescope in absolute value:
\(\sum_s|\Delta\beta_s|=F\tau\) and
\(\sum_s|\Delta\alpha_s|=F\tau^2/2m\), \textbf{independently of the
number of instruments and of their times}. An equal-prior test with
error probability at most \(\epsilon\) has
\(\mathrm{TV}(P_{\mathrm I},P_{\mathrm F})\ge1-2\epsilon\). \(\square\)

The monotonicity of \(\alpha\) and \(\beta\) is the whole of
protocol-universality. Inserting marks subdivides the displacement
without enlarging its total variation, so a protocol with a thousand
instruments faces the same budget as one with none. This is the
probabilistic counterpart of Corollary 4 of the derivation note, and it
is the point at which the present framework answers Ozawa: an instrument
whose disturbance is correlated with its outcome changes the conditional
states \(\sigma\), and therefore matters only through the apertures
\(L\) and \(P\) that those states respect.

\textbf{Corollary 3 (the floor at a given aperture).} Under Theorem 2,
with \(A=LP\) and \(r=L/\sqrt{\tau A/2m}\),

\[\tau\Delta E=\frac{F^2\tau^3}{2m}\ \ge\ \frac{2m\,\theta^2\hbar^2\tau}{(2mL+\tau P)^2}
=\frac{\theta^2\hbar^2}{A\,(r+r^{-1})^2},\qquad\theta=\arcsin(1-2\epsilon),\]

and at the balanced aperture of minimal area, \(r=1\) and \(A=\hbar/2\),
that is \(L=\sqrt{\hbar\tau/4m}\) and \(P=\sqrt{m\hbar/\tau}\),
\(\tau\Delta E\ge\theta^2\hbar/2\).

\emph{Proof.} Theorem 2 with the sharp constant gives
\(F\ge2m\hbar\theta/[\tau(2mL+\tau P)]\), so
\(\tau\Delta E\ge2m\hbar^2\theta^2\tau/(2mL+\tau P)^2\). With
\(L=r\sqrt{\tau A/2m}\) one has \(2mL=r\sqrt{2m\tau A}\) and
\(\tau P=\tau A/L=r^{-1}\sqrt{2m\tau A}\), so the denominator is
\(2m\tau A(r+r^{-1})^2\). \(\square\)

The floor is positive exactly when both sides of the aperture are
bounded, and it falls as either side grows. At fixed area it is largest
at \(r=1\), with value \(\theta^2\hbar^2/(4A)\), and it tends to zero as
\(r\to0\) or \(r\to\infty\): a minimum-uncertainty Gaussian with
\(L\to\infty\) and \(P=\hbar/2L\) keeps \(A=\hbar/2\) and distinguishes
the hypotheses at any force. So the resource the floor prices is the
aperture's two sides, and the area \(A\) controls it only once the shape
is fixed.

\hypertarget{the-counterexample-universality-fails-without-the-aperture}{%
\subsection{4. The counterexample: universality fails without the
aperture}\label{the-counterexample-universality-fails-without-the-aperture}}

\textbf{Proposition 4.} Fix \(m,F,\tau>0\) and put
\(\eta=\tau\Delta E/\hbar\). For every \(\eta<1\) there are a pure
preparation and a final measurement distinguishing the hypotheses with
error probability at most \(C\eta\) for an absolute constant \(C\). The
preparation has position spread of order \(\hbar/(F\tau)\), which
diverges as \(\eta\to0\) at fixed \(\tau\) and \(m\).

\emph{Proof.} Write \(\alpha=F\tau^2/2m\) and \(\beta=F\tau\), so
\(\alpha\beta=\tau\Delta E=\eta\hbar\). Let \(\chi\) be a fixed real
normalized profile supported in \([-1/2,1/2]\) with
\(\|\chi'\|_2<\infty\), and set

\[\psi=\frac{1}{\sqrt2}\left(\chi_++\chi_-\right),\qquad
\chi_\pm(y)=w^{-1/2}\chi\!\left(\frac{y\mp d/2}{w}\right),\qquad
d=\frac{\pi\hbar}{\beta},\qquad w=\sqrt{\frac{\alpha\hbar}{\beta}} .\]

Then \(w/d=\sqrt{\alpha\beta/\hbar}/\pi=\sqrt\eta/\pi\) and
\(\alpha/w=\sqrt{\alpha\beta/\hbar}=\sqrt\eta\), so for \(\eta<1\) the
two packets are disjoint and the displacement is small compared with the
packet width. Up to a global phase,
\((D(\alpha,\beta)\psi)(y)=e^{i\beta(y-\alpha/2)/\hbar}\psi(y-\alpha)\).
Since the supports of \(\chi_+\) and \(\chi_-\) are separated by
\(d-w>d/2\) while the shift is \(\alpha<w\), the cross terms vanish and

\[\langle\psi|D|\psi\rangle=\frac12\sum_{\pm}e^{\pm i\beta d/(2\hbar)}
\int \overline{\chi_\pm(y)}\,e^{i\beta(y\mp d/2)/\hbar}\chi_\pm(y-\alpha)\,dy .\]

With \(d=\pi\hbar/\beta\) the two prefactors are \(e^{\pm i\pi/2}\) and
cancel the leading terms, leaving the two residuals. Each residual
integral differs from \(1\) by at most
\(\|\chi(\cdot)-\chi(\cdot-\alpha/w)\|_2 +\beta w/\hbar\le(\alpha/w)\|\chi'\|_2+\beta w/\hbar\),
so

\[|\langle\psi|D|\psi\rangle|\le\frac{\alpha}{w}\|\chi'\|_2+\frac{\beta w}{\hbar}
=\sqrt\eta\left(\|\chi'\|_2+1\right).\]

Helstrom's optimal test then has error probability
\(\frac12(1-\sqrt{1-|\langle\psi|D|\psi\rangle|^2})\le\frac14|\langle\psi|D|\psi\rangle|^2 \le C\eta\)
with \(C=(\|\chi'\|_2+1)^2/4\). The position spread of \(\psi\) is
\(d/2+O(w)=\pi\hbar/(2F\tau)+O(w)\). \(\square\)

The counterexample saturates Theorem 2. Its aperture has
\(L\simeq\pi\hbar/(2F\tau)\), for which the first term of the theorem's
left side is about \(\pi/2\), so the necessary condition is met with no
room to spare. The momentum spread is of order
\(\hbar/w=\sqrt{2m\hbar/\tau}\), which stays bounded as \(\eta\to0\);
the diverging resource is the separation, not the energy. So the
obstruction to reading Newton's shrinking sagitta is a bound on
\textbf{how large the apparatus may be}, and a laboratory of unbounded
extent has no Planck gap in this comparison.

\hypertarget{position-against-the-standard-quantum-limit}{%
\subsection{5. Position against the standard quantum
limit}\label{position-against-the-standard-quantum-limit}}

The combination \(F^2\tau^3\gtrsim m\hbar\) is the standard quantum
limit for detecting a force on a free mass
(\href{https://doi.org/10.1017/CBO9780511622748}{Braginsky and Khalili,
\emph{Quantum Measurement}, 1992}, metadata). Its status has been
disputed since Yuen's objection:
\href{https://doi.org/10.1103/PhysRevLett.54.2465}{Caves defended it}
(PRL \textbf{54}, 2465, 1985, abstract) and
\href{https://doi.org/10.1103/PhysRevLett.60.385}{Ozawa exhibited a
measurement breaking it for free-mass position} (PRL \textbf{60}, 385,
1988, abstract). The present statement sits beside that dispute, and
three differences fix its place.

\begin{itemize}
\tightlist
\item
  \textbf{Universality over protocols is proved, not assumed.} The usual
  derivations fix a monitoring scheme and balance its imprecision
  against its back-action. Theorem 2 quantifies over all finite
  sequences of instruments with arbitrary adaptivity, and the reason it
  can is structural: the signal enters as a phase-space path of bounded
  total variation, and subdividing a path does not lengthen it.
\item
  \textbf{The resource is named and the bound is tight in it.} Ozawa's
  construction and Proposition 4 both buy distinguishability with
  preparations of large phase-space extent. Corollary 3 prices that
  exchange at each aperture \((L,P)\), and Proposition 4 shows the price
  is right, so the standard quantum limit appears here as the balanced
  aperture of area \(\hbar/2\) rather than as a law.
\item
  \textbf{The criterion is hypothesis testing, not estimation.} The
  quantity bounded is the error probability of deciding between two
  known histories, which is the question Newton's geometry asks, and not
  the variance of an estimate of \(F\).
\end{itemize}

Relative to the error--disturbance literature, the argument uses no such
relation. It uses only Mandelstam--Tamm and the uncertainty relation
through \(A\ge\hbar/2\), which is why the contested calibration
questions of Ozawa and of
\href{https://doi.org/10.1103/PhysRevLett.111.160405}{Busch, Lahti and
Werner} (PRL \textbf{111}, 160405, 2013, abstract) do not enter.

\hypertarget{what-this-settles-for-the-programme}{%
\subsection{6. What this settles for the
programme}\label{what-this-settles-for-the-programme}}

Against the classical theorem of the
\href{planck-gap-derivation.md}{derivation note} the parallel is exact
in form and opposite in direction. Classically the floor of the recorded
comparison is \(\tau\Delta E\ge\frac92\kappa\), with
\(\kappa=\delta\Delta\) the phase-space cost of a single mark, so
\textbf{coarser marks raise the floor}. Quantum mechanically the floor
at aperture \((L,P)\) is
\(2m\arcsin^2(1-2\epsilon)\hbar^2\tau/(2mL+\tau P)^2\), so \textbf{a
larger laboratory, in either direction of phase space, lowers it}. The
two meet at the balanced aperture of the smallest area the uncertainty
relation permits: there \(A=\hbar/2\) and the floor is \(\hbar/2\),
which is the order of the value \(\kappa=\Lambda p=h/2\) that Newton's
interval of fits and corpuscle impulse supply in the
\href{newton-mark-floor.md}{mark-floor note}.

The Newton-age reading is therefore sharper than before. Newton's limit
in Lemmas X and XI takes the sagitta and the swept area to zero, and
nothing in the geometry stops it. What stops the \emph{recorded}
comparison is the phase-space area of the apparatus, and the two
Newtonian quantities that enter are precisely the two components of the
displacement: the sagitta \(F\tau^2/2m\) against the momentum spread,
and the impulse \(F\tau\) of Proposition I against the position spread.
Newton had both quantities and had no reason to pair them with spreads
of anything.

\hypertarget{consequence-for-state}{%
\subsection{7. Consequence for STATE}\label{consequence-for-state}}

The goal of proving the probabilistic form is discharged. Theorem 1
gives the single-shot bound with the sharp constant
\(\arccos(2\sqrt{\epsilon(1-\epsilon)})\), Theorem 2 extends it to every
finite adaptive protocol of instruments, with any Kraus rank and any
apparatus memory, with the constant \(1-2\epsilon\), sharpened to
\(\arcsin(1-2\epsilon)\) by the path-length note, Corollary 3 converts
it to the floor at each aperture, \(\arcsin^2(1-2\epsilon)\hbar^2/(4A)\)
at the balanced one, and Proposition 4 decides the universality
question: universality over instruments holds, and universality over
preparations fails, with an explicit family that saturates the theorem.
Positioning against the standard quantum limit is in Section 5, and no
error--disturbance relation is used.

Open, in order:

\begin{enumerate}
\def\labelenumi{\arabic{enumi}.}
\tightlist
\item
  \textbf{Mixed conditional states: done, 2026-09-22.} The first version
  of Theorem 2 assumed rank-one Kraus operators. Purification closes the
  general case with the aperture imposed on the body's own marginal,
  since \(\Delta_\psi(G\otimes\mathbf 1)=\Delta_{\sigma_{\rm body}}G\).
\item
  \textbf{The sharp constant: done, 2026-09-23.} The Bures-angle chain
  gives Theorem 2 the constant \(\theta(\epsilon)=\arcsin(1-2\epsilon)\)
  of Theorem 1 (path-length note, Theorem P). The worst-case interval
  \([9,36]\) of the derivation note is a separate question.
\item
  \textbf{The general force law: done, 2026-09-23.} The kick form
  \(\int_0^\tau|f|L\ge\hbar\theta(\epsilon)\) of the path-length note
  holds for every force history.
\item
  \textbf{Publication.} With Section 5 in place the foundations paper
  has its positioning; what remains before submission is item 1, item 2
  and a literature pass on resource-bounded quantum metrology, where a
  bound of the shape of Corollary 3 may already exist for estimation
  rather than testing.
\end{enumerate}
\end{document}
